\documentclass[11pt]{article}
\usepackage[margin=1in]{geometry}
\usepackage[T1]{fontenc}
\usepackage{lmodern}
\usepackage{amsmath,amssymb,amsthm}
\usepackage{xcolor}
\usepackage[hidelinks]{hyperref}
\usepackage{lea-tags}

\definecolor{leagreen}{HTML}{155E50}
\newtheorem{definition}{Definition}[section]
\newtheorem{lemma}[definition]{Lemma}
\newtheorem{theorem}[definition]{Theorem}
\newtheorem{corollary}[definition]{Corollary}
\newtheorem{claim}[definition]{Claim}
\newcommand{\N}{\mathbb{N}}
\newcommand{\R}{\mathbb{R}}
\newcommand{\D}{D}
\newcommand{\Q}{Q}
\newcommand{\demotitle}[2]{%
  \noindent{\small\color{leagreen}\sffamily\bfseries #1}\par
  \vspace{0.3em}
  {\LARGE\sffamily\bfseries #2}\par
  \vspace{0.8em}}
\setlength{\parindent}{0pt}
\setlength{\parskip}{0.5em}
\hypersetup{pdftitle={From a mathematical claim to checked Lean},
  pdfauthor={Lea demo}}
\begin{document}

% Upload this whole project ZIP to Overleaf. Keep main.tex as the main document.
% The presenter script and copy/paste snippets live OUTSIDE this project.
% Do not upload them as extra .tex files: Lea inventories the whole project.
% The source deliberately contains one false claim, visibly identified in the PDF.
% Initial inventory: 9 marked items. Adding the live marker makes 10.

\demotitle{LEA DEMO / 01}{A tiny mathematical project}
\section{Doubling natural numbers}

We start with a definition, prove a structural fact, and use it in a
corollary. The mathematics is deliberately small so we can focus on the
formalization workflow.

\begin{definition}[Doubling]\label{def:demo-double}
\leadefinition{
  label=demo_double,
  context={Use a computational definition named demo_double with explicit argument (n : Nat), result Nat, and body n + n. The notation D in the source denotes this definition.}
}
For a natural number $n\in\mathbb{N}$, define $D(n)=n+n$.
\end{definition}

\begin{lemma}[Doubling zero]\label{lem:demo-double-zero}
% DEMO LIVE INSERT: add the marker from the presenter script below this line.
% This lemma is intentionally UNMARKED in the starting project.
The double of zero is zero: $D(0)=0$.
\end{lemma}
\begin{proof}
By the definition of $D$, we have $D(0)=0+0=0$.
\end{proof}

% A custom environment is intentional: its inline tag supplies the target kind.
\begin{claim}[Doubling a sum]\label{clm:demo-double-add}
\leatheorem{
  label=demo_double_add,
  uses={demo_double},
  context={Use the existing demo_double definition. Name the theorem demo_double_add and keep both arguments explicit in the order (m n : Nat). Follow the author's rearrangement proof.}
}
For every $m,n\in\mathbb{N}$,
\[
  D(m+n)=D(m)+D(n).
\]
\end{claim}
\begin{proof}
% lea: proof-for=demo_double_add
Expand the definition and rearrange addition:
\[
  (m+n)+(m+n)=(m+m)+(n+n).
\]
\end{proof}

\begin{corollary}[Adding two before doubling]\label{cor:demo-double-two}
\leacorollary{
  label=demo_double_two,
  uses={demo_double,demo_double_add},
  context={Use the existing definition and apply demo_double_add n 2 with explicit arguments. Name this theorem demo_double_two with explicit argument (n : Nat). Retain that lemma application in the proof instead of proving the whole result independently.}
}
For every $n\in\mathbb{N}$,
\[
  D(n+2)=D(n)+4.
\]
\end{corollary}
\begin{proof}
% lea: proof-for=demo_double_two
Apply the preceding claim with $m=n$ and with its second argument equal
to $2$. Then use $D(2)=2+2=4$.
\end{proof}

\vfill
{\small\color{leagreen}Definition $\longrightarrow$ structural fact
$\longrightarrow$ corollary. The source document remains ordinary LaTeX.}
\clearpage
\demotitle{LEA DEMO / 02}{When the source needs attention}
\section{A deliberately flawed draft}

This exercise starts with an intentionally false claim and an invalid proof.
The presenter notes give the replacement text to use after inspecting
Lea's finding.

\begin{theorem}[Square versus the original number]\label{thm:demo-square-bound}
% lea: formalize label=demo_square_bound
% lea: context={Read this current statement and its associated proof literally over Real. Do not silently add a hypothesis, weaken the conclusion, or treat the factored expression as nonnegative without justification. If there is a source obstruction, publish the finding and request author input under the enabled source-pause policy. A checked counterexample is also a valid diagnosis of the stated claim.}
For every real number $x\in\mathbb{R}$,
\[
  x^2\ge x.
\]
\end{theorem}
\begin{proof}
% lea: proof-for=demo_square_bound
We have $x^2-x=x(x-1)$. Both factors are nonnegative for every real
number $x$, so their product is nonnegative. Therefore $x^2\ge x$.
\end{proof}

The aim is to inspect the reported issue, clarify the mathematical
statement, and bring the formal artifact back into agreement with the
revised source.

\vfill
{\small\color{leagreen}Read the statement, the author's proof, the Lean
artifact, and the scope of the assessment together.}

\clearpage
\demotitle{LEA DEMO / 03}{One supporting result, two consumers}
\section{A shifted square}

\begin{definition}[Shifted square]\label{def:demo-shifted-square}
% lea: define label=demo_shifted_square
% lea: context={Use a computational definition named demo_shifted_square with explicit arguments (a x : Real) in that order, result Real, and body (x - a) ^ 2. The notation Q_a(x) denotes this definition.}
For $a,x\in\mathbb{R}$, set $Q_a(x)=(x-a)^2$.
\end{definition}

% Intentionally unmarked: this is supporting source mathematics, not an item.
% Do not add its LaTeX label to uses. The consumers request a shared helper.
\begin{lemma}[Shift identity]\label{lem:demo-shift-identity}
For all $a,t\in\mathbb{R}$,
\[
  Q_a(a+t)=t^2.
\]
\end{lemma}
\begin{proof}
Substitute the definition and simplify $(a+t)-a=t$.
\end{proof}

\begin{theorem}[One unit from the center]\label{thm:demo-shift-one}
% lea: formalize label=demo_shift_one uses={demo_shifted_square}
% lea: context={Follow the author's proof using the unmarked Shift identity. Search the project library and Mathlib first. If this exact result for the custom definition is absent, call request_project_dependency with mode deferred, name demo_shift_identity, kind theorem, and signature (a t : Real) : demo_shifted_square a (a + t) = t ^ 2 using the existing definition's canonical name and import. Deferred mode is intentional for this sharing demonstration even though the proof is short. Use the host-assigned module and declaration. Do not replace the shared result with a private helper or directly expand this consumer.}
For every $a\in\mathbb{R}$, $Q_a(a+1)=1$.
\end{theorem}
\begin{proof}
% lea: proof-for=demo_shift_one
Apply the Shift identity with $t=1$, then simplify $1^2=1$.
\end{proof}

\begin{theorem}[Two units from the center]\label{thm:demo-shift-two}
% lea: formalize label=demo_shift_two uses={demo_shifted_square}
% lea: context={Apply the same unrestricted Shift identity as the first consumer. Search for the checked shared helper demo_shift_identity and register reuse through request_project_dependency with kind theorem and signature (a t : Real) : demo_shifted_square a (a + t) = t ^ 2 using the existing definition's canonical name and import. Request mode deferred only if it is not already available. Use the host-assigned import and declaration. Do not reprove the identity independently or reserve a specialized version.}
For every $a\in\mathbb{R}$, $Q_a(a+2)=4$.
\end{theorem}
\begin{proof}
% lea: proof-for=demo_shift_two
Apply the same Shift identity with $t=2$, then simplify $2^2=4$.
\end{proof}

\vfill
{\small\color{leagreen}The unmarked identity supplies source context.
The shared-results view shows whether a formal helper was actually created
and reused.}

\clearpage
\demotitle{LEA DEMO / 04}{Two optional input forms}
\section{Standalone prose and existing Lean}

The main demo uses comments and tags inside environments. These two
small extras show that an enclosing theorem environment is optional and
that existing Lean can also be supplied as the source material.

\subsection{A standalone proposition tag}
\leaproposition{
  label=demo_double_one,
  uses={demo_double},
  context={Use the existing demo_double definition to show its value at one is two.}
}{The double of one is two: $D(1)=2$.}

\subsection{A Lean code block}
The following is a statement with a complete proof. The code is the input
to Lea; its presence in the PDF alone is not evidence of a checked Lea run.

\begin{leacode}{label=demo_nat_zero, context={Preserve this Nat statement and its proof.}}
theorem demo_nat_zero (n : Nat) : n + 0 = n := by
  rfl
\end{leacode}

\vfill
{\small\color{leagreen}These examples are optional in a 15-minute
recording. They can also supply small independent items for batch work.}

\end{document}
